\documentclass[11pt,a4paper]{article}

\usepackage[dutch]{babel}
\usepackage[T1]{fontenc}
\usepackage[utf8]{inputenc}

\usepackage{amsmath,amssymb,amsfonts}
\usepackage{physics}
\usepackage{mathtools}
\usepackage{siunitx}
\usepackage{geometry}
\usepackage{hyperref}
\usepackage{enumitem}
\usepackage{tcolorbox}
\usepackage{graphicx}

\geometry{margin=2.3cm}
\setlist[itemize]{topsep=4pt,itemsep=2pt}
\setlist[enumerate]{topsep=4pt,itemsep=6pt}

\title{De nabla (del) operator \texorpdfstring{$\nabla$}{∇}: werking, betekenis en toepassingen in de fysica}
\author{Jeroen Vermeulen}
\date{}

\begin{document}
    \maketitle

    \begin{tcolorbox}[colback=gray!6,colframe=gray!40,title=Doel van dit document]
        Dit document legt de nabla-operator \(\nabla\) diepgaand uit, met nadruk op:
        \begin{itemize}
            \item \textbf{Wat \(\nabla\) precies is} (een differentiaaloperator) en hoe ze werkt op verschillende objecten.
            \item \textbf{De drie kernbewerkingen} die eruit voortkomen: gradi\"ent, divergentie en rotatie (curl).
            \item \textbf{Fysische interpretaties} (flux, bronsterkte, werveling) en waar dit terugkomt in o.a. elektromagnetisme, vloeistoffen en golfvoortplanting.
            \item \textbf{Rekenen in verschillende co\"ordinaten} (Cartesisch, cilinder, bol).
            \item \textbf{Oefeningen met uitgewerkte oplossingen}.
        \end{itemize}
    \end{tcolorbox}

    \tableofcontents
    \newpage

% ============================================================


    \section{Wat is de nabla-operator?}
    De nabla-operator (ook \emph{del}-operator) is een \textbf{vector-differentiaaloperator}. In Cartesische co\"ordinaten \((x,y,z)\) defini\"eren we
    \begin{equation}
        \nabla \;\coloneqq\; \vb*{e}_x \pdv{}{x} + \vb*{e}_y \pdv{}{y} + \vb*{e}_z \pdv{}{z}.
    \end{equation}
    Hier zijn \(\vb*{e}_x,\vb*{e}_y,\vb*{e}_z\) de eenheidsvectoren, en \(\pdv{}{x}\) enz. partiële afgeleiden.

    \subsection{Belangrijk: \(\nabla\) is geen gewone vector}
    Hoewel \(\nabla\) er uitziet als een vector, is het \textbf{een operator} die pas betekenis krijgt wanneer ze op een veld werkt.

    Je kan haar zien als:
    \begin{itemize}
        \item een \textbf{richting-gevoelige differentiatormachine};
        \item een compact notatiesysteem dat meerdere afgeleiden bundelt;
        \item een object dat zich formeel als een vector gedraagt onder co\"ordinatentransformaties, maar met differentiaaloperatoren als componenten.
    \end{itemize}

    \subsection{Waarop kan \(\nabla\) werken?}
    \(\nabla\) kan werken op:
    \begin{itemize}
        \item een \textbf{scalair veld} \(\phi(\vb{r})\) (bv. temperatuur, potentiaal, dichtheid);
        \item een \textbf{vectorveld} \(\vb{A}(\vb{r})\) (bv. snelheidsveld, elektrisch veld);
        \item producten/combinaties van velden (waar dan productregels, kettingregels, enz. spelen).
    \end{itemize}

% ============================================================


    \section{De drie kernoperators uit \texorpdfstring{$\nabla$}{∇}}
    Uit \(\nabla\) volgen drie fundamentele differentiaaloperatoren:
    \begin{enumerate}
        \item \textbf{Gradi\"ent:} \(\nabla \phi\)
        \item \textbf{Divergentie:} \(\nabla\cdot\vb{A}\)
        \item \textbf{Rotatie (curl):} \(\nabla\times\vb{A}\)
    \end{enumerate}
    En daarnaast: de \textbf{Laplaciaan} \(\nabla^2\), die je kan zien als \(\nabla\cdot\nabla\) toegepast op een scalair veld.

% ------------------------------------------------------------

    \subsection{Gradi\"ent \texorpdfstring{$\nabla \phi$}{∇φ}}
    Laat \(\phi(x,y,z)\) een scalair veld zijn. Dan
    \begin{equation}
        \nabla \phi =
        \vb*{e}_x \pdv{\phi}{x} + \vb*{e}_y \pdv{\phi}{y} + \vb*{e}_z \pdv{\phi}{z}.
    \end{equation}

    \subsubsection{Geometrische betekenis}
    \(\nabla \phi\) wijst in de richting van de \textbf{steilste toename} van \(\phi\).
    De grootte \(\abs{\nabla \phi}\) is de \textbf{maximale stijgsnelheid per lengte-eenheid}.

    Formeel: voor een kleine verplaatsing \(d\vb{r}\) geldt
    \begin{equation}
        d\phi = \nabla\phi \cdot d\vb{r}.
    \end{equation}
    Dit is het \textbf{totale differentiaal} in vectorvorm: de gradi\"ent is de unieke vector die deze relatie voor alle \(d\vb{r}\) vervult.

    \subsubsection{Niveaumanifolds en orthogonaliteit}
    Op een niveauoppervlak \(\phi(\vb{r})=\text{const}\) is \(d\phi=0\), dus
    \begin{equation}
        0 = \nabla\phi\cdot d\vb{r}.
    \end{equation}
    Dat betekent: \(\nabla\phi\) staat loodrecht op alle verplaatsingen \(d\vb{r}\) die in het niveauoppervlak liggen. \\
    \(\Rightarrow\) \(\nabla\phi\) is \textbf{normaal} op het niveauoppervlak.

    \subsubsection{Fysische voorbeelden}
    \begin{itemize}
        \item Elektrostatica: \(\vb{E} = -\nabla V\) (elektrisch veld is minus gradi\"ent van potentiaal).
        \item Mechanica: \(\vb{F} = -\nabla U\) (kracht uit potenti\"ele energie).
        \item Diffusie: flux \(\vb{J}\) is vaak proportioneel aan \(-\nabla c\) (Fick).
        \item Warmtegeleiding: \(\vb{q} = -k\nabla T\) (Fourier).
    \end{itemize}

% ------------------------------------------------------------

    \subsection{Divergentie \texorpdfstring{$\nabla\cdot\vb{A}$}{∇·A}}
    Laat \(\vb{A}(x,y,z) = (A_x,A_y,A_z)\). Dan
    \begin{equation}
        \nabla\cdot\vb{A} = \pdv{A_x}{x}+\pdv{A_y}{y}+\pdv{A_z}{z}.
    \end{equation}

    \subsubsection{Lokale bronsterkte: flux per volume}
    Divergentie meet de \textbf{netto uitstroming} (flux naar buiten) per volume-eenheid in het limietgeval van een heel klein volume.

    Heuristisch: neem een klein volume \(V\) rond punt \(\vb{r}\), met gesloten oppervlak \(S=\partial V\). Dan:
    \begin{equation}
        \nabla\cdot\vb{A}(\vb{r}) = \lim_{V\to 0}\frac{1}{V}\iint_{\partial V}\vb{A}\cdot d\vb{S}.
    \end{equation}
    Dit is de \textbf{differenti\"ele vorm} van de Gauss-divergentiestelling.

    \subsubsection{Fysische interpretatie}
    \begin{itemize}
        \item \(\nabla\cdot\vb{A} > 0\): punt gedraagt zich als \textbf{bron} (meer uitstroom dan instroom).
        \item \(\nabla\cdot\vb{A} < 0\): \textbf{put} (sink).
        \item \(\nabla\cdot\vb{A} = 0\): \textbf{soleno\"idaal} veld (incompressibel/bronvrij in die regio).
    \end{itemize}

    \subsubsection{Voorbeelden in fysica}
    \begin{itemize}
        \item Elektromagnetisme: \(\nabla\cdot\vb{E}=\rho/\varepsilon_0\) (ladingdichtheid is bron van \(\vb{E}\)).
        \item Magnetisme: \(\nabla\cdot\vb{B}=0\) (geen magnetische monopolen).
        \item Vloeistoffen: continu\"iteit \(\partial_t\rho+\nabla\cdot(\rho\vb{v})=0\). In incompressibele stroming (\(\rho\) constant): \(\nabla\cdot\vb{v}=0\).
    \end{itemize}

% ------------------------------------------------------------

    \subsection{Rotatie (curl) \texorpdfstring{$\nabla\times\vb{A}$}{∇×A}}
    In Cartesische co\"ordinaten:
    \begin{equation}
        \nabla\times\vb{A} =
        \begin{vmatrix}
            \vb*{e}_x & \vb*{e}_y & \vb*{e}_z\\
            \pdv{}{x} & \pdv{}{y} & \pdv{}{z}\\
            A_x & A_y & A_z
        \end{vmatrix}.
    \end{equation}

    \subsubsection{Lokale circulatie: wervelsterkte}
    Curl meet de \textbf{circulatie per oppervlakte} in het limietgeval:
    \begin{equation}
        (\nabla\times\vb{A})\cdot \vb{n} = \lim_{S\to 0}\frac{1}{S}\oint_{\partial S}\vb{A}\cdot d\vb{l},
    \end{equation}
    waar \(S\) een klein oppervlakje is met normaal \(\vb{n}\), en \(\partial S\) de randcurve.

    Dit is de \textbf{differenti\"ele vorm} van de Stokes-stelling.

    \subsubsection{Fysische interpretatie}
    \begin{itemize}
        \item Curl nul: veld is \textbf{irrotationeel} (lokaal geen werveling).
        \item Curl niet nul: er is lokale \textbf{rotatie/werveling}.
    \end{itemize}

    \subsubsection{Voorbeelden}
    \begin{itemize}
        \item Elektrostatica: \(\nabla\times\vb{E}=0\) (conservatief veld).
        \item Faraday: \(\nabla\times\vb{E}=-\partial_t \vb{B}\).
        \item Stromingsleer: vorticiteit \(\boldsymbol{\omega}=\nabla\times\vb{v}\).
    \end{itemize}

% ------------------------------------------------------------

    \subsection{Laplaciaan \texorpdfstring{$\nabla^2$}{∇²}}
    De Laplaciaan op een scalair veld \(\phi\) is:
    \begin{equation}
        \nabla^2\phi \;=\; \nabla\cdot(\nabla\phi)
        \;=\; \pdv[2]{\phi}{x}+\pdv[2]{\phi}{y}+\pdv[2]{\phi}{z}.
    \end{equation}

    \subsubsection{Interpretatie (intu\"itief)}
    \(\nabla^2\phi\) meet in zekere zin \textbf{hoeveel \(\phi\) lokaal afwijkt van haar gemiddelde} in een kleine omgeving.
    In veel fysica verschijnt dit als \textbf{diffusie/spreiding/curvature-term}.

    \subsubsection{Typische fysische vergelijkingen}
    \begin{itemize}
        \item Poisson: \(\nabla^2 V = -\rho/\varepsilon_0\).
        \item Laplace: \(\nabla^2 V=0\) (bronvrije regio).
        \item Warmtevergelijking: \(\partial_t T = \alpha \nabla^2 T\).
        \item Golfvergelijking: \(\partial_t^2 \psi = c^2 \nabla^2 \psi\).
        \item Schr\"odinger: \(-\frac{\hbar^2}{2m}\nabla^2\psi + V\psi = i\hbar \partial_t\psi\).
    \end{itemize}

% ============================================================


    \section{Productregels en nuttige vectoridentiteiten}
    Omdat \(\nabla\) een differentiaaloperator is, gelden productregels.
    Laat \(\phi\) scalair en \(\vb{A},\vb{B}\) vectorvelden zijn.

    \subsection{Productregels}
    \begin{align}
        \nabla(\phi\psi) &= \phi\nabla\psi + \psi\nabla\phi,\\
        \nabla\cdot(\phi\vb{A}) &= \nabla\phi\cdot\vb{A} + \phi(\nabla\cdot\vb{A}),\\
        \nabla\times(\phi\vb{A}) &= \nabla\phi\times\vb{A} + \phi(\nabla\times\vb{A}).
    \end{align}

    \subsection{Identiteiten (zeer belangrijk in fysica)}
    \begin{align}
        \nabla\times(\nabla\phi) &= \vb{0} \qquad \text{(curl van gradi\"ent is nul)},\\
        \nabla\cdot(\nabla\times\vb{A}) &= 0 \qquad \text{(div van curl is nul)},\\
        \nabla\times(\nabla\times\vb{A}) &= \nabla(\nabla\cdot\vb{A}) - \nabla^2 \vb{A}.
    \end{align}
    Die laatste is cruciaal om golfvergelijkingen voor EM-velden af te leiden.

% ============================================================


    \section{Co\"ordinatenstelsels: waarom \texorpdfstring{$\nabla$}{∇} ingewikkelder wordt}
    In Cartesische co\"ordinaten zijn basisvectoren constant.
    In kromlijnige co\"ordinaten (cilindrisch, bol) hangen basisvectoren en lengteschalen af van de positie.
    Daarom verschijnen \textbf{schaalfactoren}.

    \subsection{Cilindrische co\"ordinaten \((r,\varphi,z)\)}
    \begin{align}
        \nabla \phi &=
        \vb*{e}_r \pdv{\phi}{r} +
        \vb*{e}_\varphi \frac{1}{r}\pdv{\phi}{\varphi} +
        \vb*{e}_z \pdv{\phi}{z},\\[4pt]
        \nabla\cdot\vb{A} &=
        \frac{1}{r}\pdv{}{r}\!\big(rA_r\big) + \frac{1}{r}\pdv{A_\varphi}{\varphi} + \pdv{A_z}{z},\\[4pt]
        \nabla^2\phi &=
        \frac{1}{r}\pdv{}{r}\!\left(r\pdv{\phi}{r}\right)
        + \frac{1}{r^2}\pdv[2]{\phi}{\varphi}
        + \pdv[2]{\phi}{z}.
    \end{align}

    \subsection{Bolco\"ordinaten \((r,\theta,\varphi)\)}
    \begin{align}
        \nabla \phi &=
        \vb*{e}_r \pdv{\phi}{r} +
        \vb*{e}_\theta \frac{1}{r}\pdv{\phi}{\theta} +
        \vb*{e}_\varphi \frac{1}{r\sin\theta}\pdv{\phi}{\varphi},\\[4pt]
        \nabla^2\phi &=
        \frac{1}{r^2}\pdv{}{r}\!\left(r^2\pdv{\phi}{r}\right)
        +\frac{1}{r^2\sin\theta}\pdv{}{\theta}\!\left(\sin\theta\pdv{\phi}{\theta}\right)
        +\frac{1}{r^2\sin^2\theta}\pdv[2]{\phi}{\varphi}.
    \end{align}

    \begin{tcolorbox}[colback=yellow!7,colframe=yellow!40,title=Praktische tip]
        In fysica kies je co\"ordinaten die de symmetrie volgen:
        \begin{itemize}
            \item puntlading / sferische massa \(\rightarrow\) bolco\"ordinaten;
            \item oneindige draad / cilinder \(\rightarrow\) cilindrisch;
            \item vlakke problemen \(\rightarrow\) Cartesisch.
        \end{itemize}
        Dan worden \(\nabla\)-bewerkingen vaak drastisch eenvoudiger.
    \end{tcolorbox}

% ============================================================


    \section{Waar komt \texorpdfstring{$\nabla$}{∇} overal terug in de fysica?}
    Hier is het kernidee: \(\nabla\) vertaalt \textbf{lokale veranderingen} naar \textbf{globale wetten} via integraalstellingen.

    \subsection{Elektrostatica en potentiaaltheorie}
    \begin{align}
        \vb{E} &= -\nabla V,\\
        \nabla\cdot\vb{E} &= \rho/\varepsilon_0 \quad \Rightarrow \quad \nabla^2 V = -\rho/\varepsilon_0.
    \end{align}
    Interpretatie: ladingdichtheid \(\rho\) is bron van veldlijnen; potentiaal voldoet aan Poisson.

    \subsection{Maxwellvergelijkingen (differentiaalvorm)}
    \begin{align}
        \nabla\cdot\vb{E} &= \rho/\varepsilon_0,\\
        \nabla\cdot\vb{B} &= 0,\\
        \nabla\times\vb{E} &= -\pdv{\vb{B}}{t},\\
        \nabla\times\vb{B} &= \mu_0\vb{J}+\mu_0\varepsilon_0\pdv{\vb{E}}{t}.
    \end{align}
    Met vectoridentiteiten kan je hieruit EM-golfvergelijkingen afleiden.

    \subsection{Continu\"iteit en behoudswetten}
    \begin{equation}
        \pdv{\rho}{t} + \nabla\cdot(\rho \vb{v}) = 0.
    \end{equation}
    Dit is massabehoud: verandering in dichtheid in een volume = netto instroom.

    \subsection{Vloeistofdynamica (Navier--Stokes, schets)}
    Voor een Newtoniaanse vloeistof:
    \begin{quote}
        Een \textbf{Newtoniaanse vloeistof} is een vloeistof waarbij de schuifspanning lineair evenredig is met de schuifsnelheid:
        \(\tau = \mu\,\dot\gamma\). De viscositeit \(\mu\) is (bij gegeven temperatuur en druk) ongeveer constant en hangt dus niet af van de vervormingssnelheid zelf.
        Voorbeelden zijn water, lucht en (bij benadering) lichte oli\"en.
    \end{quote}
    \begin{equation}
        \rho\left(\pdv{\vb{v}}{t}+(\vb{v}\cdot\nabla)\vb{v}\right)
        = -\nabla p + \mu \nabla^2 \vb{v} + \rho \vb{f}.
    \end{equation}
    Hier zie je:
    \begin{itemize}
        \item \(-\nabla p\): kracht door drukgradi\"ent;
        \item \(\mu\nabla^2\vb{v}\): viscositeit/diffusie van impuls;
        \item \((\vb{v}\cdot\nabla)\vb{v}\): convectieterm (niet-lineair).
    \end{itemize}

    \subsection{Diffusie en warmte}
    \begin{equation}
        \pdv{u}{t} = D\nabla^2 u.
    \end{equation}
    \(\nabla^2\) is de operator die ruimtelijke spreiding encodeert.

    \subsection{Golven}
    \begin{equation}
        \pdv[2]{\psi}{t} = c^2 \nabla^2 \psi.
    \end{equation}
    In 3D: de Laplaciaan bepaalt hoe kromming van \(\psi\) leidt tot versnelling in de tijd.

% ============================================================


    \section{Oefeningen (met volledige uitwerkingen)}

    \subsection*{Notatie}
    We werken standaard in Cartesische co\"ordinaten tenzij anders vermeld.
    Gebruik: \(\vb{r}=(x,y,z)\).

% ------------------------------------------------------------

    \subsection{Oefening 1: gradi\"ent en richting van snelste stijging}
    Gegeven \(\phi(x,y,z)=x^2y+3z\).
    \begin{enumerate}
        \item Bereken \(\nabla\phi\).
        \item Bepaal \(\nabla\phi\) in het punt \((1,2,-1)\).
        \item Geef de eenheidsrichting van snelste stijging in dat punt.
    \end{enumerate}

    \subsubsection*{Oplossing}
    \begin{enumerate}
        \item
        \[
            \pdv{\phi}{x}=2xy,\qquad \pdv{\phi}{y}=x^2,\qquad \pdv{\phi}{z}=3.
        \]
        Dus
        \[
            \nabla\phi=(2xy)\vb*{e}_x+(x^2)\vb*{e}_y+3\vb*{e}_z.
        \]
        \item In \((1,2,-1)\):
        \[
            \nabla\phi(1,2,-1) = (2\cdot1\cdot2)\vb*{e}_x + (1^2)\vb*{e}_y + 3\vb*{e}_z
            = 4\vb*{e}_x+\vb*{e}_y+3\vb*{e}_z.
        \]
        \item De richting van snelste stijging is de genormaliseerde gradi\"ent:
        \[
            \hat{\vb{u}} = \frac{\nabla\phi}{\abs{\nabla\phi}}
            = \frac{(4,1,3)}{\sqrt{4^2+1^2+3^2}}
            = \frac{(4,1,3)}{\sqrt{26}}.
        \]
    \end{enumerate}

% ------------------------------------------------------------

    \subsection{Oefening 2: divergentie en bron/put interpretatie}
    Beschouw \(\vb{A}(x,y,z)=(x,\,y,\,z)\).
    \begin{enumerate}
        \item Bereken \(\nabla\cdot\vb{A}\).
        \item Interpreteer fysisch (bron/put?) in elk punt.
        \item Bereken de flux door de bol \(x^2+y^2+z^2=R^2\) en controleer met Gauss.
    \end{enumerate}

    \subsubsection*{Oplossing}
    \begin{enumerate}
        \item
        \[
            \nabla\cdot\vb{A}=\pdv{x}{x}+\pdv{y}{y}+\pdv{z}{z}=1+1+1=3.
        \]
        \item Divergentie is overal positief en constant: elk punt gedraagt zich als een \emph{uniforme bron}.
        \item Op de bol is \(\vb{A}=\vb{r}\). De uitwendige normaal is \(\hat{\vb{n}}=\vb{r}/R\).
        Dus op het oppervlak:
        \[
            \vb{A}\cdot \hat{\vb{n}} = \vb{r}\cdot \frac{\vb{r}}{R}=\frac{r^2}{R}=\frac{R^2}{R}=R.
        \]
        Flux:
        \[
            \Phi = \iint \vb{A}\cdot d\vb{S} = \iint (\vb{A}\cdot\hat{\vb{n}})\, dS = \iint R\, dS = R \cdot 4\pi R^2 = 4\pi R^3.
        \]
        Gauss-controle:
        \[
            \Phi = \iiint_V (\nabla\cdot \vb{A})\, dV = \iiint_V 3\, dV = 3\cdot \frac{4}{3}\pi R^3 = 4\pi R^3.
        \]
        Consistent.
    \end{enumerate}

% ------------------------------------------------------------

    \subsection{Oefening 3: curl en conservativiteit}
    Gegeven \(\vb{F}(x,y,z)=(-y,\,x,\,0)\).
    \begin{enumerate}
        \item Bereken \(\nabla\times \vb{F}\).
        \item Is \(\vb{F}\) conservatief in \(\mathbb{R}^3\)? Motiveer.
        \item Bereken de circulatie \(\oint_C \vb{F}\cdot d\vb{l}\) rond de cirkel \(x^2+y^2=R^2\) in het vlak \(z=0\) (tegenwijzerzin).
    \end{enumerate}

    \subsubsection*{Oplossing}
    \begin{enumerate}
        \item
        \[
            \nabla\times\vb{F}=
            \begin{vmatrix}
                \vb*{e}_x & \vb*{e}_y & \vb*{e}_z\\
                \pdv{}{x} & \pdv{}{y} & \pdv{}{z}\\
                -y & x & 0
            \end{vmatrix}
            =
            \left(0-0\right)\vb*{e}_x
            -\left(0-0\right)\vb*{e}_y
            +\left(\pdv{x}{x}-\pdv{(-y)}{y}\right)\vb*{e}_z
            =
            (1-(-1))\vb*{e}_z = 2\vb*{e}_z.
        \]
        \item Omdat \(\nabla\times\vb{F}\neq \vb{0}\), is \(\vb{F}\) niet irrotationeel en dus niet conservatief (globaal) in \(\mathbb{R}^3\).
        \item Parameteriseer \(C\): \(x=R\cos t,\; y=R\sin t,\; t\in[0,2\pi]\).
        Dan \(d\vb{l}=(\dot x,\dot y,0)\,dt=(-R\sin t, R\cos t,0)\,dt\).
        En \(\vb{F}=(-y,x,0)=(-R\sin t, R\cos t,0)\).
        Dus:
        \[
            \vb{F}\cdot d\vb{l} = (-R\sin t, R\cos t,0)\cdot(-R\sin t, R\cos t,0)\,dt = R^2(\sin^2 t+\cos^2 t)\,dt=R^2\,dt.
        \]
        Circulatie:
        \[
            \oint_C \vb{F}\cdot d\vb{l} = \int_0^{2\pi} R^2\,dt = 2\pi R^2.
        \]
    \end{enumerate}

% ------------------------------------------------------------

    \subsection{Oefening 4: Laplaciaan en harmonische functies}
    Laat \(\phi(x,y,z)=x^2+y^2-2z^2\).
    \begin{enumerate}
        \item Bereken \(\nabla^2\phi\).
        \item Is \(\phi\) harmonisch? (d.w.z. voldoet \(\nabla^2\phi=0\)?)
    \end{enumerate}

    \subsubsection*{Oplossing}
    \[
        \pdv[2]{\phi}{x}=2,\qquad \pdv[2]{\phi}{y}=2,\qquad \pdv[2]{\phi}{z}=-4.
    \]
    Dus
    \[
        \nabla^2\phi = 2+2-4=0.
    \]
    Ja, \(\phi\) is harmonisch.

% ------------------------------------------------------------

    \subsection{Oefening 5: van integraal naar differentiaal (continu\"iteit)}
    Stel dat een hoeveelheid \(Q\) met dichtheid \(\rho(\vb{r},t)\) stroomt met flux \(\vb{J}(\vb{r},t)\).
    Je weet dat voor elk vast volume \(V\):
    \[
        \dv{}{t}\iiint_V \rho\, dV = - \iint_{\partial V} \vb{J}\cdot d\vb{S}.
    \]
    \begin{enumerate}
        \item Gebruik de divergentietheorema om de rechterkant om te schrijven als een volume-integraal.
        \item Leid af dat lokaal geldt: \(\partial_t\rho + \nabla\cdot\vb{J}=0\).
    \end{enumerate}

    \subsubsection*{Oplossing}
    \begin{enumerate}
        \item Divergentiestelling:
        \[
            \iint_{\partial V}\vb{J}\cdot d\vb{S} = \iiint_V (\nabla\cdot\vb{J})\, dV.
        \]
        Dus:
        \[
            \dv{}{t}\iiint_V \rho\, dV = -\iiint_V (\nabla\cdot\vb{J})\, dV.
        \]
        \item Breng alles samen:
        \[
            \iiint_V \pdv{\rho}{t}\, dV + \iiint_V (\nabla\cdot\vb{J})\, dV = 0
            \quad\Rightarrow\quad
            \iiint_V \left(\pdv{\rho}{t} + \nabla\cdot\vb{J}\right)dV = 0.
        \]
        Omdat dit voor \emph{elk} volume \(V\) geldt, moet de integrand overal nul zijn:
        \[
            \pdv{\rho}{t} + \nabla\cdot\vb{J}=0.
        \]
    \end{enumerate}

% ------------------------------------------------------------

    \subsection{Oefening 6: overgang van Cartesische naar bolcoördinaten}

    Beschouw het scalair veld

    \[
        \phi(x,y,z) = x^2 + y^2 + z^2.
    \]

    \begin{enumerate}
        \item Schrijf de Cartesische coördinaten \((x,y,z)\) uit in functie van bolcoördinaten \((r,\theta,\varphi)\).

        \item Druk vervolgens het veld \(\phi(x,y,z)\) uit in bolcoördinaten.

        \item Bereken de gradiënt \(\nabla \phi\) in Cartesische coördinaten.

        \item Bereken de gradiënt \(\nabla \phi\) rechtstreeks in bolcoördinaten met de formule

        \[
            \nabla \phi =
            \vb{e}_r \pdv{\phi}{r} +
            \vb{e}_\theta \frac{1}{r}\pdv{\phi}{\theta} +
            \vb{e}_\varphi \frac{1}{r\sin\theta}\pdv{\phi}{\varphi}.
        \]

        \item Toon aan dat beide resultaten fysisch equivalent zijn.
    \end{enumerate}


% ------------------------------------------------------------

    \subsection*{Oplossing}

    \begin{enumerate}

        \item De relatie tussen Cartesische en bolcoördinaten is

        \[
            x = r\sin\theta\cos\varphi
        \]

        \[
            y = r\sin\theta\sin\varphi
        \]

        \[
            z = r\cos\theta
        \]

        waar

        \begin{itemize}
            \item \(r\) de afstand tot de oorsprong is,
            \item \(\theta\) de poolhoek (hoek t.o.v. de \(z\)-as),
            \item \(\varphi\) de azimutale hoek in het \(xy\)-vlak.
        \end{itemize}

        \item Substitueer in \(\phi(x,y,z)\):

        \[
            \phi = x^2 + y^2 + z^2
        \]

        \[
            = (r^2\sin^2\theta\cos^2\varphi)
            + (r^2\sin^2\theta\sin^2\varphi)
            + (r^2\cos^2\theta)
        \]

        Factoriseer \(r^2\):

        \[
            \phi = r^2(\sin^2\theta\cos^2\varphi
            + \sin^2\theta\sin^2\varphi
            + \cos^2\theta)
        \]

        Gebruik

        \[
            \cos^2\varphi + \sin^2\varphi = 1
        \]

        dus

        \[
            \phi = r^2(\sin^2\theta + \cos^2\theta)
        \]

        \[
            \phi = r^2
        \]

        Het scalair veld hangt dus enkel af van \(r\).

        \item Gradiënt in Cartesische coördinaten:

        \[
            \nabla\phi =
            \left(
                \pdv{\phi}{x},
                \pdv{\phi}{y},
                \pdv{\phi}{z}
            \right)
        \]

        \[
            \pdv{\phi}{x} = 2x, \quad
            \pdv{\phi}{y} = 2y, \quad
            \pdv{\phi}{z} = 2z
        \]

        Dus

        \[
            \nabla\phi = (2x,2y,2z)
        \]

        Of in vectorvorm

        \[
            \nabla\phi = 2(x,y,z) = 2\vb{r}.
        \]

        \item In bolcoördinaten hebben we

        \[
            \phi = r^2
        \]

        dus

        \begin{gather*}
            \pdv{\phi}{r} = 2r\\
            \pdv{\phi}{\theta} = 0\\
            \pdv{\phi}{\varphi} = 0\\
        \end{gather*}

        Invullen in de gradiëntformule:

        \[
            \nabla\phi =
            \vb{e}_r (2r)
        \]

        Dus

        \[
            \nabla\phi = 2r\,\vb{e}_r.
        \]

        \item Controle van de equivalentie

        De radiale eenheidsvector is

        \[
            \vb{e}_r = \frac{\vb{r}}{r}.
        \]

        Dus

        \begin{gather*}
            2r\,\vb{e}_r =
            2r \frac{\vb{r}}{r}\\
            = 2\vb{r}\\
        \end{gather*}

        Dit is exact hetzelfde resultaat als in Cartesische coördinaten:

        \[
            \nabla\phi = 2(x,y,z).
        \]

        De resultaten zijn dus volledig consistent.

    \end{enumerate}


    \section{De "Nabla-piramide" (Hiërarchie van Velden)}
    Het is voor beginners vaak verwarrend welk type veld er in een operator gaat en wat eruit komt. Een kort overzichtje helpt:
    \begin{center}
        \begin{tabular}{c|c|c}
            Operator         & Input                  & Output                                           \\
            \hline
            \(\nabla\)       & scalair veld \(\phi\)  & vector veld \(\nabla\phi\) (gradi\"ent)          \\
            \(\nabla\cdot\)  & vector veld \(\vb{A}\) & scalair veld \(\nabla\cdot\vb{A}\) (divergentie) \\
            \(\nabla\times\) & vector veld \(\vb{A}\) & vector veld \(\nabla\times\vb{A}\) (curl)        \\
            \(\nabla^2\)     & scalair veld \(\phi\)  & scalair veld \(\nabla^2\phi\) (Laplaciaan)
        \end{tabular}
    \end{center}

    \subsection{Uitgewerkt voorbeeld bij elke operator uit de tabel}
    \begin{tcolorbox}[colback=blue!4,colframe=blue!35,title=Van type veld naar concrete berekening]
        \textbf{1) Gradi\"ent \(\nabla\phi\): van scalair veld naar vectorveld}\\
        Neem
        \[
            \phi(x,y,z)=x^2y+2yz-z^2.
        \]
        Dan
        \[
            \pdv{\phi}{x}=2xy,\qquad
            \pdv{\phi}{y}=x^2+2z,\qquad
            \pdv{\phi}{z}=2y-2z.
        \]
        Dus
        \[
            \nabla\phi=(2xy)\vb*{e}_x+(x^2+2z)\vb*{e}_y+(2y-2z)\vb*{e}_z.
        \]
        In het punt \((1,-1,2)\):
        \[
            \nabla\phi(1,-1,2)=(-2,5,-6).
        \]
        \emph{Interpretatie:} deze vector wijst in de richting van de snelste lokale toename van \(\phi\).

        \medskip
        \textbf{2) Divergentie \(\nabla\cdot\vb{A}\): van vectorveld naar scalair veld}\\
        Neem
        \[
            \vb{A}(x,y,z)=\big(x^2,\,xy,\,-2z\big).
        \]
        Dan
        \[
            \nabla\cdot\vb{A}=\pdv{x^2}{x}+\pdv{xy}{y}+\pdv{(-2z)}{z}=2x+x-2=3x-2.
        \]
        In het vlak \(x=1\):
        \[
            \nabla\cdot\vb{A}=1>0.
        \]
        \emph{Interpretatie:} lokaal gedraagt het veld zich daar als een bron (netto uitstroom).

        \medskip
        \textbf{3) Curl \(\nabla\times\vb{A}\): van vectorveld naar vectorveld}\\
        Neem
        \[
            \vb{A}(x,y,z)=(-y,\,x,\,0).
        \]
        Dan
        \[
            \nabla\times\vb{A}=
            \begin{vmatrix}
                \vb*{e}_x & \vb*{e}_y & \vb*{e}_z\\
                \pdv{}{x} & \pdv{}{y} & \pdv{}{z}\\
                -y & x & 0
            \end{vmatrix}
            =\left(0-0\right)\vb*{e}_x-\left(0-0\right)\vb*{e}_y
            +\left(\pdv{x}{x}-\pdv{(-y)}{y}\right)\vb*{e}_z
            =2\vb*{e}_z.
        \]
        \emph{Interpretatie:} de werveling is constant en wijst overal langs de positieve \(z\)-as.

        \medskip
        \textbf{4) Laplaciaan \(\nabla^2\phi\): van scalair veld naar scalair veld}\\
        Neem
        \[
            \phi(x,y,z)=x^2+y^2-2z^2.
        \]
        Dan
        \[
            \nabla^2\phi=\pdv[2]{\phi}{x}+\pdv[2]{\phi}{y}+\pdv[2]{\phi}{z}=2+2-4=0.
        \]
        \emph{Interpretatie:} \(\phi\) is harmonisch; er is lokaal geen netto bronterm in dit veld.
    \end{tcolorbox}


    \section{De Helmholtz-decompositien}
    Dit is een fundamenteel concept in de fysica dat direct aansluit op hoofdstuk 3.2 over identiteiten\textsuperscript{}.
    Het stelt dat (onder milde voorwaarden) elk vectorveld $A$ geschreven kan worden als de som van een irrotationeel deel (een gradiënt) en een solenoïdaal deel (een rotatie):
    \[
        \vb{A} = \nabla\phi + \nabla\times\vb{B},
    \]
    Dit verklaart \textit{waarom} we in de fysica zoveel belang hechten aan
    $\nabla \cdot A = 0$ (bronvrij) en
    $\nabla \times A = 0$
    (conservatief)\textsuperscript{}\textsuperscript{}\textsuperscript{}\textsuperscript{}.

% ============================================================


    \section{Samenvatting in 10 regels}
    \begin{tcolorbox}[colback=green!6,colframe=green!35,title=Cheat sheet]
        \begin{itemize}
            \item \(\nabla = (\partial_x,\partial_y,\partial_z)\) (Cartesisch): vector-differentiaaloperator.
            \item \(\nabla\phi\): gradi\"ent = richting en grootte van snelste stijging.
            \item \(\nabla\cdot\vb{A}\): divergentie = netto flux naar buiten per volume (bron/put).
            \item \(\nabla\times\vb{A}\): curl = circulatie/werveling per oppervlak.
            \item \(\nabla^2\phi = \nabla\cdot\nabla\phi\): Laplaciaan = diffusie/curvature-operator.
            \item Gauss: \(\iiint (\nabla\cdot\vb{A})dV = \iint \vb{A}\cdot d\vb{S}\).
            \item Stokes: \(\iint (\nabla\times\vb{A})\cdot d\vb{S} = \oint \vb{A}\cdot d\vb{l}\).
            \item \(\nabla\times(\nabla\phi)=0\), \(\nabla\cdot(\nabla\times\vb{A})=0\).
            \item Maxwell, continu\"iteit, Navier--Stokes, diffusie en golven zijn doordrenkt van \(\nabla\).
            \item Kies co\"ordinaten volgens symmetrie: dan wordt \(\nabla\) hanteerbaar.
        \end{itemize}
    \end{tcolorbox}

\end{document}


